\section*{अध्याय \ref{ch:b3:seismology-imaging} --- भूकंप विज्ञान और पृथ्वी का प्रतिबिंबन}

\begin{solution}{exo:b3:seismology-imaging:1}
(क) $\sin i_0 = pv_0/R = 500\times5.80/6371 = 0.4552$, अतः $i_0 = \ang{27.1}$।
(ख) \qty{771}{km} पर, $r = \qty{5600}{km}$ और $r/v = 5600/11.066 =
\qty{506.1}{s.rad^{-1}}$; \qty{871}{km} पर, $r = \qty{5500}{km}$ और $r/v =
5500/11.245 = \qty{489.1}{s.rad^{-1}}$। जहाँ $r/v = p = 500$ होता है, वहाँ किरण पलटती है:
$z_t = 771 + 100\times(506.1 - 500)/(506.1 - 489.1) = 771 + 36 =
\qty{807}{km}$।
\end{solution}

\begin{solution}{exo:b3:seismology-imaging:2}
(क) $(455.5 - 369.6)/10 = \qty{8.59}{s}$ प्रति अंश, और
$8.59\times180/\pi = \qty{492.2}{s.rad^{-1}}$। (ख) चाप का एक अंश
$6371\pi/180 = \qty{111.19}{km}$ है, अतः तरंग-अग्र सतह को
$111.19/8.59 = \qty{12.94}{km.s^{-1}}$ की दर से पार करता है, जो सतह की किसी भी चट्टान से कहीं अधिक तेज़ है;
अतः किरण तीखे कोण से पहुँचती है। (ग) \omterm{def:b3:seismology-imaging:ray-parameter}{पलटाव बिंदु} पर $v = r_t/p =
5450/492.2 = \qty{11.07}{km.s^{-1}}$, निचले प्रावार में P तरंग की चाल।
\end{solution}

\begin{solution}{exo:b3:seismology-imaging:3}
(क) $\mu_1 = \rho_1\beta_1^2 = 2700\times3500^2 = \qty{3.31e10}{Pa}$
(\qty{33.1}{GPa}); $\mu_2 = 3300\times4500^2 = \qty{6.68e10}{Pa}$
(\qty{66.8}{GPa})। (ख) $\lambda = cT = 3.79\times20 = \qty{75.8}{km}$, परत की मोटाई का 2.17
गुना: तरंग अर्ध-अंतरिक्ष में काफ़ी गहराई तक पहुँचती है, जो
इसी कारण इसकी चाल पहले ही $\beta_1$ से ऊपर है। (ग) $4000/3.384 = \qty{1182}{s}$ और
$4000/4.333 = \qty{923}{s}$: 20-s तरंगें 100-s तरंगों से \qty{259}{s}, या
\qty{4.3}{min} बाद पहुँचती हैं।
\end{solution}

\begin{solution}{exo:b3:seismology-imaging:4}
(क) ${}_0S_2$: $1/(53.89\times60) = \qty{3.093e-4}{Hz} = \qty{0.3093}{mHz}$;
${}_0T_2$: $1/(44.01\times60) = \qty{0.3787}{mHz}$। (ख) $1440/53.89 = 26.7$
और $1440/44.01 = 32.7$ दोलन प्रतिदिन। (ग) अंतर
$\qty{0.0694}{mHz}$ है, जिसका व्युत्क्रम \qty{14400}{s} है: कम से कम \qty{4.0}{h} का
अभिलेख चाहिए, व्यवहार में इससे कहीं अधिक, क्योंकि क्षीणन से शिखर चौड़े हो जाते हैं और
विभाजित हो जाते हैं।
\end{solution}

\begin{solution}{exo:b3:seismology-imaging:5}
(क) गोले के दो बिंदुओं के बीच की जीवा, जो केंद्र पर कोण $\Delta$ से अलग है,
इसकी लंबाई $2R\sin(\Delta/2)$ है, और इसे $v$ चाल से तय किया जाता है।
(ख) \ang{30} पर: $2\times6371\times\sin\ang{15}/10.0 = \qty{329.8}{s}$, जबकि प्रेक्षित मान
\qty{369.6}{s} है; \ang{90} पर: $2\times6371\times\sin\ang{45}/10.0 =
\qty{901.0}{s}$, जबकि प्रेक्षित मान \qty{779.7}{s} है। (ग) \ang{30} वाली किरण पलटती है
ऊपरी प्रावार में, जहाँ वास्तविक चाल लगभग 8 से \qty{10}{km.s^{-1}} है, जो
एकसमान \qty{10}{km.s^{-1}} से कम है; \ang{90} वाली किरण अपना अधिकांश पथ
निचले प्रावार में बिताती है, जहाँ चाल 11 से \qty{13.7}{km.s^{-1}} है, और उसमें नीचे की ओर झुकती है, इसलिए
वास्तविक पृथ्वी अधिक तेज़ है।
\end{solution}

\begin{solution}{exo:b3:seismology-imaging:6}
(क) वर्ग करने पर, $(2 - \xi)^4 = 16(1 - \xi/3)(1 - \xi)$। विस्तार करने पर:
$16 - 32\xi + 24\xi^2 - 8\xi^3 + \xi^4 = 16 - \tfrac{64}{3}\xi + \tfrac{16}{3}\xi^2$,
अतः $\xi^4 - 8\xi^3 + \tfrac{56}{3}\xi^2 - \tfrac{32}{3}\xi = 0$, अर्थात
$\xi\,(3\xi^3 - 24\xi^2 + 56\xi - 32) = 0$, और
$3\xi^3 - 24\xi^2 + 56\xi - 32 = (\xi - 4)(3\xi^2 - 12\xi + 8)$। मूल
$\xi = 0$ कोई तरंग नहीं है। (ख) द्विघात के मूल $2 \pm 2/\sqrt3$ हैं,
अर्थात 3.155 और 0.8453; केवल $\xi = 0.8453$ अंतराल $(0, 1)$ में आता है (और $\xi = 4$
नहीं आता), और यह अवर्गित समीकरण को संतुष्ट करता है, दोनों पक्ष $4/3$ के बराबर हैं।
अतः $c/\beta = \sqrt{0.8453} = 0.9194$। (ग) $c = 0.9194\times3.46 =
\qty{3.181}{km.s^{-1}}$; \qty{1000}{km} की दूरी पर, $314.4 - 289.0 = \qty{25.3}{s}$
S तरंग के बाद।
\end{solution}

\begin{solution}{exo:b3:seismology-imaging:7}
$\mathsf G$ की पंक्तियाँ $(1, x_i)$ हैं और $\mathsf C_d^{-1} = \operatorname{diag}
(100, 100, 25, 25)$ (\unit{s^{-2}})। प्रसामान्य समीकरण इस प्रकार हैं
\[
  \begin{pmatrix} 250 & 23750 \\ 23750 & \num{2812500} \end{pmatrix}
  \begin{pmatrix} a \\ b \end{pmatrix}
  = \begin{pmatrix} 3665 \\ 418125 \end{pmatrix},
\]
जिसका निर्धारक \num{1.390625e8} है। अतः $b = (250\times418125 -
23750\times3665)/\num{1.390625e8} = \qty{0.12575}{s.km^{-1}}$ और
$a = (2812500\times3665 - 23750\times418125)/\num{1.390625e8} =
\qty{2.713}{s}$। सहप्रसरण $\mathsf N^{-1}$ है: $\sigma_b =
\sqrt{250/\num{1.390625e8}} = \qty{0.00134}{s.km^{-1}}$, $\sigma_a =
\sqrt{2812500/\num{1.390625e8}} = \qty{0.142}{s}$। चाल
$1/b = \qty{7.95}{km.s^{-1}}$ है, जिसकी त्रुटि $\sigma_b/b^2 =
\qty{0.08}{km.s^{-1}}$: सबसे ऊपरी प्रावार की चाल। अवशिष्ट हैं $-0.001$,
$0.011$, $-0.076$ और $\qty{0.036}{s}$, और $\chi^2 = 0.19$, $4 - 2 = 2$
स्वातंत्र्य-कोटि के लिए: अपेक्षा से कम, अतः बताई गई त्रुटियाँ, यदि कुछ हो,
तो बड़ी हैं; समंजन स्वीकार्य है।
\end{solution}

\begin{solution}{exo:b3:seismology-imaging:8}
(क) $\mathsf G = (30\;\;40)$ एक ही पंक्ति है: $\mathsf G^{\mathsf T}\mathsf G$
कोटि एक का $2\times2$ आव्यूह है, अतः वह अव्युत्क्रमणीय है; $30m_1 + 40m_2 = 0$ वाला कोई भी मॉडल
जोड़ा जा सकता है, बिना आँकड़े को बदले। (ख) $\mathsf G\mathsf G^{\mathsf T} =
900 + 1600 = 2500$, अतः $\mathbf m = (30, 40)\times0.50/2500 =
(0.006, 0.008)\,\unit{s.km^{-1}}$। (ग) $\mathsf R = \frac{1}{2500}
\begin{pmatrix} 900 & 1200 \\ 1200 & 1600 \end{pmatrix} =
\begin{pmatrix} 0.36 & 0.48 \\ 0.48 & 0.64 \end{pmatrix}$, ट्रेस 1 का: एक
आँकड़ा प्राचलों के एक संयोजन को विभेदित करता है। (घ) $\mathsf R\,(0.01, 0) =
(0.0036, 0.0048)\,\unit{s.km^{-1}}$: विसंगति किरण के अनुदिश फैलती है, अधिकांश
उस कोष्ठ में जाती है जहाँ किरण लंबी है।
\end{solution}

\begin{solution}{exo:b3:seismology-imaging:9}
(क) रैखिकृत: $\delta t = -(\delta v/v)L/v = -0.03\times200/9.0 =
\qty{-0.667}{s}$; सटीक रूप से, $L/(1.03v) - L/v = 200/9.27 - 200/9.0 =
\qty{-0.647}{s}$। (ख) स्लैब में पथ $100/\cos\ang{40} =
\qty{130.5}{km}$: $-0.03\times130.5/9.0 = \qty{-0.435}{s}$ रैखिकृत और
$\qty{-0.422}{s}$ सटीक। (ग) सटीक मान रैखिकृत मान को
$1.03$ से भाग देने पर मिलता है: \qty{3}{\percent} की त्रुटि, जो स्वयं $\delta v/v$ की कोटि की है, जैसा
प्रथम-कोटि के सिद्धांत से अपेक्षित है।
\end{solution}

\begin{solution}{exo:b3:seismology-imaging:10}
(क) अधिस्वर तब तक विद्यमान रहता है जब तक उसका अर्ध-अंतरिक्ष भाग क्षय होता है, $c < \beta_2$; अपनी कट-ऑफ़ पर
$c = \beta_2$, अतः $q_2 = 0$ और $\tan(kHq_1) = 0$, जहाँ $kHq_1$
$[n\pi, n\pi + \pi/2)$ में है: $kHq_1 = n\pi$। (ख) $k = \omega/\beta_2$ लेने पर और
$q_1 = \sqrt{\beta_2^2/\beta_1^2 - 1} = \sqrt{\beta_2^2 - \beta_1^2}/\beta_1$,
$\omega_nH\sqrt{\beta_2^2 - \beta_1^2}/(\beta_1\beta_2) = n\pi$, जो
सूत्र है। (ग) $\sqrt{4.5^2 - 3.5^2} = \sqrt8 = 2.828$; $\omega_1 =
\pi\times3.5\times4.5/(35\times2.828) = \qty{0.4998}{rad.s^{-1}}$, जिसका आवर्तकाल
$2\pi/0.4998 = \qty{12.57}{s}$: अधिक आवर्तकालों पर केवल मूल विधा
विद्यमान है। (घ) $n = 0$ के लिए कट-ऑफ़ आवृत्ति शून्य है: मूल विधा
सभी आवर्तकालों पर विद्यमान है। लघु आवर्तकालों पर $kH \to \infty$ जबकि $kHq_1$
$\pi/2$ से नीचे रहता है, अतः $q_1 \to 0$ और $c \to \beta_1$।
\end{solution}

\begin{solution}{exo:b3:seismology-imaging:11}
(क) $\mathsf G^{\mathsf T}(\mathsf G\mathsf G^{\mathsf T} + \varepsilon^2\mathsf I)
= (\mathsf G^{\mathsf T}\mathsf G + \varepsilon^2\mathsf I)\mathsf G^{\mathsf T}$;
बाईं ओर से $(\mathsf G^{\mathsf T}\mathsf G + \varepsilon^2
\mathsf I)^{-1}$ और दाईं ओर से $(\mathsf G\mathsf G^{\mathsf T} +
\varepsilon^2\mathsf I)^{-1}$ से गुणा कीजिए। $\varepsilon > 0$ के लिए दोनों व्युत्क्रमणीय हैं।
(ख) यहाँ $\mathsf G\mathsf G^{\mathsf T} = \|\mathbf g\|^2 = 2500$ एक संख्या है,
इसलिए $\mathsf R_\varepsilon = \mathsf G^{\mathsf T}\mathsf G/(2500 + \varepsilon^2)$,
जिसका ट्रेस $2500/(2500 + \varepsilon^2)$ है। (ग) अनुमानित आँकड़ा
$\mathbf g\cdot\mathbf m_\varepsilon = 0.50\times2500/(2500 + \varepsilon^2)$
और अवशिष्ट $0.50\,\varepsilon^2/(2500 + \varepsilon^2)$ है। मान इस प्रकार हैं:
$\varepsilon^2 = 0$ के लिए: ट्रेस 1, अवशिष्ट 0; 2500 के लिए: ट्रेस 0.5, अवशिष्ट
\qty{0.25}{s}; \num{10000} के लिए: ट्रेस 0.2, अवशिष्ट \qty{0.40}{s}।
अधिक अवमंदन से मॉडल छोटा होता है, विभेदन घटता है और बेमेल बढ़ता है।
\end{solution}

\begin{solution}{exo:b3:seismology-imaging:12}
(क) अभ्यास 5 से, $T = (2R/v)\sin(\Delta/2)$, इसलिए $\mathrm dT/\mathrm d\Delta
= (R/v)\cos(\Delta/2)$। केंद्र पर कोण $\Delta$ बनाने वाली जीवा
शीर्ष-कोण $\Delta$ वाले समद्विबाहु त्रिभुज का आधार है, इसलिए वह स्रोत पर त्रिज्या से
$i_0 = (\pi - \Delta)/2$ कोण बनाती है, और $p = R\sin i_0/v =
(R/v)\cos(\Delta/2)$: वही। (ख) \omterm{def:b3:seismology-imaging:ray-parameter}{पलटाव बिंदु} जीवा के मध्य में है,
केंद्र से $R\cos(\Delta_1/2)$ दूरी पर; समतुल्य रूप से
$r_1 = p_1v$। (ग) यहाँ $p(\Delta)/p_1 = \cos(\Delta/2)/\cos(\Delta_1/2)$, और
सूत्र से $r_1 = R\exp\bigl(-\ln(1/\cos(\Delta_1/2))\bigr) =
R\cos(\Delta_1/2)$ मिलता है, जैसा (ख) में है; चाल $r_1/p_1 = R\cos(\Delta_1/2)\big/
\bigl((R/v)\cos(\Delta_1/2)\bigr) = v$ है।
\end{solution}

\begin{solution}{pb:b3:seismology-imaging:1}
\textbf{1.} $i_0 = \arcsin 0.8 = \ang{53.13}$ और $p = 0.8/5.0 =
\qty{0.16}{s.km^{-1}}$।

\textbf{2.} किरण वहाँ मुड़ती है जहाँ $v = 1/p = \qty{6.25}{km.s^{-1}}$ है, अर्थात
$z = (6.25 - 5.0)/0.05 = \qty{25}{km}$ पर।

\textbf{3.} किरण के अनुदिश $\sin i = pv$ का अवकलन करने पर $\cos i\,\mathrm di/
\mathrm ds = pg\,\mathrm dz/\mathrm ds = pg\cos i$, अतः $\mathrm di/\mathrm ds =
pg$: दिशा स्थिर दर से घूमती है, जैसे त्रिज्या
$1/(pg) = 1/(0.16\times0.05) = \qty{125}{km}$ वाले वृत्त की वक्रता हो। \omterm{def:b3:seismology-imaging:ray-parameter}{पलटाव बिंदु} पर किरण
\qty{25}{km} गहराई पर क्षैतिज है और केंद्र ठीक \qty{125}{km}
ऊपर है: $z = \qty{-100}{km}$ पर, अर्थात सतह से $v_0/g$ ऊपर, जहाँ
बहिर्वेशित चाल शून्य हो जाती।

\textbf{4.} स्रोत पर वृत्त की त्रिज्या किरण के लंबवत है,
इसलिए वह क्षैतिज से $i_0$ कोण बनाती है, और केंद्र
क्षैतिजतः $125\cos i_0$ दूरी पर है; सममिति से
$X = 2\times125\times\cos i_0 = 2\times125\times0.6 = \qty{150}{km}$, या
$X = 2\sqrt{1 - p^2v_0^2}/(pg)$।

\textbf{5.} $\mathrm dT = \mathrm ds/v$ में $\mathrm ds = \mathrm di/(pg)$ और
$v = \sin i/p$ रखने पर: $\mathrm dT = \mathrm di/(g\sin i)$। तब
\[
  T = \frac2g\int_{i_0}^{\pi/2}\frac{\mathrm di}{\sin i} = \frac2g\ln\cot\frac{i_0}2 =
  \frac2g\operatorname{arccosh}\frac1{pv_0} .
\]
यहाँ $\cot(\ang{26.57}) = 2$, इसलिए
$T = 40\ln2 = \qty{27.73}{s}$।

\textbf{6.} $p = 0.15$ के लिए: $pv_0 = 0.75$, $\sqrt{1 - 0.5625} = 0.6614$,
$X = 2\times0.6614/0.0075 = \qty{176.4}{km}$ और $T = 40\operatorname{arccosh}
(1.3333) = 40\times0.7954 = \qty{31.81}{s}$। तब $\Delta T/\Delta X =
4.08/26.4 = \qty{0.155}{s.km^{-1}}$, 0.15 और 0.16 के बीच, जैसा
$\mathrm dT/\mathrm dX = p$ अपेक्षित है।

\textbf{7.} $\mu_1 = 2700\times3500^2 = \qty{33.1}{GPa}$ और
$\mu_2 = 3300\times4500^2 = \qty{66.8}{GPa}$।

\textbf{8.} $k = 2\pi/(cT) = 2\pi/(3.99\times25) = \qty{0.06299}{rad.km^{-1}}$;
$q_1 = \sqrt{(3.99/3.5)^2 - 1} = 0.5474$; $q_2 = \sqrt{1 - (3.99/4.5)^2} = 0.4624$।

\textbf{9.} $\tan(kHq_1) = 66.825\times0.4624/(33.075\times0.5474) = 1.707$;
मूल विधा $kHq_1 = \arctan1.707 = 1.041$ लेती है, अतः
$H = 1.041/(0.06299\times0.5474) = \qty{30.2}{km}$।

\textbf{10.} \qty{40}{s} पर: $k = 2\pi/170 = 0.03696$, $q_1 = 0.6888$,
$q_2 = 0.3287$, $\tan(kHq_1) = 0.964$, $kHq_1 = 0.7671$ और
$H = 0.7671/(0.03696\times0.6888) = \qty{30.1}{km}$। दोनों परिणाम
कला-वेगों की परिशुद्धता के भीतर मेल खाते हैं: भूपर्पटी लगभग \qty{30}{km} मोटी है।

\textbf{11.} $\omega = 0.2513$ और $\qty{0.1571}{rad.s^{-1}}$, $k = 0.06299$ और
$\qty{0.03696}{rad.km^{-1}}$: $U \approx 0.0942/0.02603 = \qty{3.62}{km.s^{-1}}$,
दोनों कला-वेगों से कम है, जैसा होना चाहिए जब $c$ आवर्त के साथ
बढ़ता है।

\textbf{12.} $\lambda = 3.99\times25 = \qty{99.8}{km}$ और $4.25\times40 =
\qty{170}{km}$। मूल विधा नीचे अपने तरंगदैर्घ्य के एक अंश तक नमूना लेती है:
25-s तरंग अधिकतर भूपर्पटी में रहती है, 40-s तरंग प्रावार में काफ़ी गहराई तक पहुँचती है,
इसलिए उसकी चाल $\beta_2$ के अधिक निकट है। दो आवर्त जो एक ही $H$ तय करते हैं,
यह जाँच है कि दो-परत मॉडल पर्याप्त है।

\textbf{13.} कोशिकाओं 1 से 4 के स्तंभों और पंक्तियों A से F के साथ, इकाई के रूप में
\qty{1000}{km} लेते हुए ($\sqrt2\times100/1000 = 0.1414$):
\[
  \mathsf G = \begin{pmatrix}
  0.1 & 0.1 & 0 & 0 \\ 0 & 0 & 0.1 & 0.1 \\ 0.1 & 0 & 0.1 & 0 \\
  0 & 0.1 & 0 & 0.1 \\ 0.1414 & 0 & 0 & 0.1414 \\ 0 & 0.1414 & 0.1414 & 0
  \end{pmatrix}.
\]

\textbf{14.} $\mathsf C_d^{-1} = \operatorname{diag}(400, 400, 400, 400, 100,
100)$। पंक्ति या स्तंभ किरण $0.1^2\times400 = 4$ का योगदान देती है, विकर्ण किरण
$0.1414^2\times100 = 2$ का योगदान देती है। कोशिका 1 को किरणें A, C (प्रत्येक 4) और E (2) पार करती हैं: 10। कोशिकाएँ
1 और 2 किरण A साझा करती हैं: 4; कोशिकाएँ 1 और 3 किरण C साझा करती हैं: 4; कोशिकाएँ 1 और 4 किरण E साझा करती हैं: 2; और
इसी प्रकार सममिति से बाकी, जिससे दी गई चार पंक्तियाँ मिलती हैं।

\textbf{15.} पंक्ति या स्तंभ किरण अपनी प्रत्येक कोशिका में $0.1\times400\,\delta t = 40\,
\delta t$ का योगदान देती है, विकर्ण किरण $0.1414\times100\,\delta t =
14.14\,\delta t$ का योगदान देती है। कोशिका 1: $40(0.16 - 0.36) + 14.14\times(-0.15) = -10.12$; कोशिका
2: $40(0.16 + 0.11) + 14.14\times(-0.27) = 6.98$; कोशिका 3: $40(-0.48 - 0.36) +
14.14\times(-0.27) = -37.42$; कोशिका 4: $40(-0.48 + 0.11) + 14.14\times(-0.15) =
-16.92$।

\textbf{16.} $(1,1,1,1)$ के लिए $\mathsf N$ की प्रत्येक पंक्ति का योग 20 है। इसी प्रकार
$(1,-1,-1,1)$ के लिए पंक्ति 1 $10 - 4 - 4 + 2 = 4$ देती है और बाकी पंक्तियाँ भी
उसी प्रतिरूप की हैं। $(1,1,-1,-1)$ के लिए $10 + 4 - 4 - 2 = 8$, और $(1,-1,1,-1)$ के लिए,
$10 - 4 + 4 - 2 = 8$ है। ये सदिश लांबिक हैं और इनका मानक 2 है, इसलिए
$\mathsf N^{-1} = \sum_k \mathbf u_k\mathbf u_k^{\mathsf T}/(4\lambda_k)$; इसका
$(1,1)$ अवयव $\tfrac14(\tfrac1{20} + \tfrac14 + \tfrac18 + \tfrac18) =
\tfrac{11}{80}$, इसका $(1,2)$ अवयव $\tfrac14(\tfrac1{20} - \tfrac14 + \tfrac18
- \tfrac18) = -\tfrac4{80}$, इसका $(1,4)$ अवयव $\tfrac14(\tfrac1{20} +
\tfrac14 - \tfrac18 - \tfrac18) = \tfrac1{80}$ है, और इसी प्रकार आगे।

\textbf{17.} $\hat{\mathbf m} = \mathsf N^{-1}\mathbf b$:
$\hat m_1 = (11\times(-10.121) - 4\times6.982 - 4\times(-37.418) - 16.921)/80
= -0.081$; इसी प्रकार $\hat m_2 = 1.844$, $\hat m_3 = -3.706$ और
$\hat m_4 = \qty{-0.931}{ms.km^{-1}}$ प्राप्त होते हैं।

\textbf{18.} $\delta v/v_0 = -v_0\,\delta s = -8.0\times10^{-3}\,\hat m$, प्रतिशत में
$-0.8\,\hat m$: कोशिका 1 $+0.07$, कोशिका 2 $-1.48$, कोशिका 3 $+2.96$, कोशिका 4
$\qty{+0.75}{\percent}$।

\textbf{19.} $\mathsf N^{-1}$ का हर विकर्ण अवयव $11/80$ है: हर कोशिका की
मानक त्रुटि $\sqrt{0.1375} = \qty{0.371}{ms.km^{-1}}$ है, अर्थात
$\qty{0.30}{\percent}$। कोशिका 1 और 2 का सहसंबंध $-4/11 = -0.36$ है:
ये किरण A साझा करती हैं, और एक कोशिका में रखी मंदता को दूसरी के विरुद्ध समायोजित किया जा सकता है।
कोशिकाएँ 1 और 4, जो केवल कम परिशुद्ध विकर्ण E साझा करती हैं, दुर्बल और
धनात्मक रूप से सहसंबद्ध हैं, $+1/11 = +0.09$।

\textbf{20.} पूर्वानुमानित अवशिष्ट $0.176$, $-0.464$, $-0.379$, $0.091$,
$-0.143$ और $\qty{-0.263}{s}$; फ़िट के बाद के अवशिष्ट $-0.016$,
$-0.016$, $0.019$, $0.019$, $-0.007$ और $\qty{-0.007}{s}$, और
$\chi^2 = 0.50$, $6 - 4 = 2$ स्वातंत्र्य कोटियों के लिए, दी गई
त्रुटियों के अनुरूप हैं, यदि कोई झुकाव है तो वह छोटे मानों की ओर है।

\textbf{21.} आइगेनसदिश $\mathsf N$ के ही हैं, और आइगेनमान
$\lambda/(\lambda + 4)$ हैं: $20/24 = 5/6$, $4/8 = 1/2$, $8/12 = 2/3$ दो बार।
प्रश्न 16 की तरह, $R_{11} = \tfrac14(\tfrac56 + \tfrac12 + \tfrac23 + \tfrac23) =
\tfrac23$; $R_{12} = \tfrac14(\tfrac56 - \tfrac12 + \tfrac23 - \tfrac23) =
\tfrac1{12}$; $R_{13} = \tfrac14(\tfrac56 - \tfrac12 - \tfrac23 + \tfrac23) =
\tfrac1{12}$; $R_{14} = \tfrac14(\tfrac56 + \tfrac12 - \tfrac23 - \tfrac23) = 0$।

\textbf{22.} $\mathbf b$ को आइगेनसदिशों पर प्रक्षेपित कीजिए: $\mathbf b\cdot\mathbf u_k/4
= -14.37$, $0.849$, $12.80$ और $-9.40$ आइगेनमान 20, 4, 8 और 8 के लिए।
अवमंदित अनुमान इन्हें $\lambda + 4$ से भाग देता है, न कि $\lambda$ से;
कोशिका 3 (घटक $1, -1, -1, 1$) के लिए,
$m_3 = -14.37/24 - 0.849/8 - 12.80/12 - 9.40/12 = \qty{-2.555}{ms.km^{-1}}$,
यह $\qty{+2.04}{\percent}$ की वेग विसंगति के बराबर है। अवमंदन हर घटक को
शून्य की ओर खींचता है, सबसे अधिक उन्हें जो कमज़ोर रूप से बाधित हैं।

\textbf{23.} अनवमंदित मॉडल, जो आँकड़ों को उनकी त्रुटियों की सीमा के भीतर फ़िट करता है,
($\chi^2 = 0.50$, 2 स्वातंत्र्य कोटियों के लिए), जबकि अवमंदित मॉडल उनसे
बुरी तरह बेमेल है ($15.8$)। अवमंदन की आवश्यकता केवल तब होती है जब $\mathsf N$ एकवचन या लगभग
ऐसा हो; यहाँ इसका सबसे छोटा आइगेनमान 4 है, जो छोटा नहीं है, अनुमान स्थिर है और
इसकी त्रुटियाँ मध्यम हैं।

\textbf{24.} पंक्ति 3 $(1/12, 0, 2/3, 1/12)$ है: कोशिका 3 का अवमंदित अनुमान
अपने वास्तविक मान के दो-तिहाई के बराबर है, साथ ही कोशिका 1 और 4 के एक-बारहवें भाग का,
किरण C और किरण B के साथ लगे उसके पड़ोसी, और कोशिका 2 का कुछ भी नहीं।

\textbf{25.} कोशिका 3, बाईं ओर की निचली कोशिका, लगभग
\qty{3}{\percent} तेज़ है; कोशिका 2 धीमी है। महासागरीय स्थलमंडल की एक ठंडी, घनी स्लैब
महाद्वीप के नीचे धँसती हुई आसपास के प्रावार से तेज़ होती है, इसलिए
गहरी और K के महासागर वाले पक्ष पर स्थित तेज़ कोशिका ही स्लैब है; उसके ऊपर की
धीमी कोशिका, दूसरी ओर, वहाँ स्थित है जहाँ प्रावार-पच्चर गर्म है।

\textbf{26.} प्रश्न 18 और 19 से, स्लैब कोशिका की P-वेग विसंगति
\textbf{\qty[text-series-to-math=true, uncertainty-mode=separate]{3.0 \pm 0.3}{\percent}} है।
\end{solution}
